\section{Introdução}

\paragraph{}
Este documento apresenta as evidências práticas das atividades desenvolvidas na
Aula 01 do módulo de Introdução ao Pentest, no âmbito do Programa Hackers do
Bem -- Nível Fundamental. A atividade tem como foco a realização de varredura
de vulnerabilidades em uma aplicação web, utilizando a ferramenta OWASP ZAP,
com o objetivo de identificar falhas de segurança com base na metodologia OWASP
Top 10 2021.

\paragraph{}
Conforme as orientações propostas, este material reúne, em um único arquivo
PDF, os registros obrigatórios da execução das atividades, incluindo: a
varredura inicial realizada com o OWASP ZAP, a análise das vulnerabilidades
identificadas, a associação dessas vulnerabilidades às categorias do OWASP Top
10 2021 e às técnicas da matriz MITRE ATT\&CK, além da elaboração de um plano
de mitigação e do relatório final consolidado. Cada etapa é acompanhada de
descrições objetivas, com o intuito de facilitar a compreensão, análise e
avaliação dos resultados obtidos.
